\documentclass[../../main.tex]{subfiles}

\graphicspath{{../../image/}{../../figure/}}

\begin{document}

\section{注释}\label{section:ch05-注释}

关于双曲守恒律方程组初值与初边值问题经典解的综合论述，见 Benzoni-Gavage 与 Serre 的专著 [2]。

对称化守恒律方程组 Cauchy 问题经典解的局部存在性已由多种方法建立，最终依赖于对系统关于空间变量微分所得的"能量"估计层级。

最早且至今仍最流行的途径阐述于 Benzoni-Gavage 与 Serre [2]，通过在\eqref{align:ch05-cauchy-system}线性化系统上的迭代过程构造解。它由 Schauder [1] 在拟线性二阶波动方程背景下首创，并经多位作者——特别地 Friedrichs [2]、Gårding [1] 与 Majda [3]——的贡献达到现今一般形式。Godunov [3]、Makino、Ukai 与 Kawashima [1]、Chemin [1]、Lax [1]、M.E. Taylor [1,2] 与 Métivier [1] 使用了凸熵 Hessian 之外的对称化子，甚至符号对称化子。就此而言，回忆 Euler 方程的对称形式\eqref{align:ch04-Euler对称双曲}，它即使在真空态也保持双曲性。

对不可压缩流体的 Euler 方程，Beale、Kato 与 Majda [1] 已证明：只要涡度保持有界，Cauchy 问题的经典解就持续存在。对可压缩流体情形，Chemin [2] 解释了经典解如何因涡度、压缩或速度散度的爆发而破坏。

\cref{theorem:ch05-凸熵局部存在}的另一建立途径由 Kato [1] 提出，基于抽象发展方程理论。此处采用粘性消去法，因为它同样适用于熵仅在对合锥上凸或仅为多凸的情形。

如\cref{theorem:ch05-相对熵稳定性}、\cref{theorem:ch05-involution-stability}与\cref{theorem:ch05-polyconvex-stability}中那样用相对熵证明经典解在更广可容许弱解类中的唯一性与稳定性（通俗称为"弱-强唯一性"），肇始于 Dafermos [9,10] 与 DiPerna [7] 的工作。该途径现已朝多个方向推广：弱解可非常弱——仅测度值——或系统可为中间双曲-抛物型，甚至出现于动理学理论。代表性样本（大量相关论文之外）为 Brenier、De Lellis 与 Székelyhidi [1]、LeFloch [8]、Christoforou 与 Tzavaras [1]、Lattanzio 与 Tzavaras [2]、Tzavaras [7]、Demoulini、Stuart 与 Tzavaras [3]、Berthelin、Tzavaras 与 Vasseur [1]、Luo 与 Smoller [1,2]、Germain [1]、Elling [5]、Miroshnikov 与 Trivisa [1]、Feireisl 与 Novotný [1]。事实上，在某些条件下该方法甚至可给出弱解的唯一性；见 DiPerna [7]、Chen、Frid 与 Li [1]、Gui-Qiang Chen 与 Yachun Li [1,2]、Gui-Qiang Chen 与 Jun Chen [1]、Kwon [3]、Kwon 与 Vasseur [2]、Choi 与 Vasseur [1]、Serre 与 Vasseur [1]、Feireisl、Kreml 与 Vasseur [1]、Leger 与 Vasseur [2]。由最后一篇文献，Texier 与 Zumbrun [2] 证明相对熵条件蕴涵任意强度极端激波的 Lopatinskii（稳定）条件。

有一些有趣的经典解例子，其中相对熵产生（\eqref{align:ch05-after-cutoff}右端第二项）恰为非正。
于是\eqref{align:ch05-relative-stability}以 $b=0$ 成立，给出一致时间的稳定性。此类条件可由状态向量的适当选取诱导。例如直线等熵气体流动中的稀疏波情形即如此。见 Gui-Qiang Chen [7]、Chen 与 Frid [7] 与 Serre [30]。

带对合的双曲守恒律方程组由 Boillat [4] 与 Dafermos [14] 讨论。特别地，Boillat [4] 给出广义相对论中出现的例子。偶然熵概念归功于 Serre [22]。\cref{section:ch05-对合与偶然熵}的分析遵循并推广 Dafermos [14,27,37]。
对合与补偿紧性理论之间的紧密关系由 Murat 与 Tartar 形式化；见 Tartar [1,2]。就此而言，"对合锥"对应"特征锥"。特别地，对弹性动力学方程，见 Hughes、Kato 与 Marsden [1] 及 Dafermos 与 Hrusa [1]。

\cref{section:ch05-偶然熵与多凸性}遵循 Dafermos [27]，它改进了本书早期版本中对该主题的处理。
如\cref{section:ch04-注记}已注，弹性静力学中多凸概念由 Ball [1] 引入，
作为使内能函数弱下半连续的条件。作者最初从 P.G. LeFloch 处听到通过附加伸长张量不变量守恒律来延拓弹性动力学守恒律方程组的想法。显式延拓首先由 Qin [1] 及 Demoulini、Stuart 与 Tzavaras [2] 发表。另见 Lattanzio 与 Tzavaras [1]。Tzavaras [8] 在松弛格式中嵌入弹性动力学延拓系统。特别地，他讨论气体动力学方程的例子，其在 Euler 公式中装备凸熵，而 Lagrangian 公式中的熵仅多凸。Brenier [2] 给出电动力学方程带 Born--Infeld 本构关系的两种不同延拓，包括此处记录的一种，并讨论其在各种 regime 下的渐近行为。该研究在 Brenier [4] 与 Brenier 与 Yong [1] 中继续。另见 Neves 与 Serre [1]。Serre [22] 利用偶然熵--熵通量对\eqref{align:ch03-BI-PL}，
在一般本构关系下设计电动力学中的恰当延拓。

关于满足 Kawashima 型条件的带源平衡律方程组整体定义经典解的存在性与长时间行为，文献广泛。该条件的变体见于 Kawashima 及其合作者的多篇论文，但其首次出现于 Shizuta 与 Kawashima [1]，故常称为\emph{Kawashima--Shizuta 条件}。读者可在 Hanouzet 与 Natalini [1]、Yong [6]、Yang、Zhu 与 Zhao [3]、Bianchini、Hanouzet 与 Natalini [1]、Kawashima 与 Yong [1,2]、Xu 与 Kawashima [1,2]、Yanni Zeng [6] 中找到\cref{theorem:ch05-damping-global}风味的结果，以及这些论文的参考文献中。另见 Beauchard 与 Zuazua [1]、Hu 与 Wang [1]、Luo、Xin 与 Zeng [1]、Mascia 与 Natalini [1]、Peng 与 Wang [1]。关于粘性耗散作用的综述见 Tai-Ping Liu [30]。

\cref{section:ch05-阻尼与松弛的作用}中一般松弛框架的设定遵循 Chen、Levermore 与 Liu [1]。
有趣的另一框架见 Bouchut [1]。松弛与扩散之间的联系最早在气体动理学理论中被认识，其中通过 Chapman--Enskog 展开实现（如 Cercignani [1]）。Chapman--Enskog 型展开也被用于联系\eqref{align:ch05-relaxed-system}类双曲平衡律方程组与\eqref{align:ch04-黏性组}形式带扩散系统；见 Kawashima 与 Yong [1,2]。

关于松弛理论各方面的文献汗牛充栋。历史回顾见 Mascia [4]。综述与广泛文献见于 Natalini [3] 与 Yong [4]。相关参考文献包括 Tai-Ping Liu [21]、Nishibata 与 Yu [1]、Wei-Cheng Wang 与 Zhouping Xin [1]、Donatelli 与 Marcati [1]、Hsiao 与 Pan [1]、Shen 与 Winther [1]、Yong [2,3,5]、Yang 与 Zhu [1]、Yang、Zhu 与 Zhao [3]、Liu 与 Yong [1]、Natalini 与 Terracina [1]、Xin 与 Xu [1]、DiFrancesco 与 Lattanzio [1]、Fan 与 Härterich [1]、Fan 与 Luo [1]、Bedaoui、Klingenberg 与 LeFloch [1]、Berthelin 与 Bouchut [1]、Junca 与 Rascle [1]、Tadmor 与 Tang [2]、Carbou 与 Hanouzet [1]、Carbou、Hanouzet 与 Natalini [1]、Chalons 与 Coulombel [1]、Lambert 与 Marchesin [1]、Yanni Zeng [4,5]、Miroshnikov 与 Trivisa [2]、Lattanzio 与 Tzavaras [1]。特别地，Jin 与 Xin [1] 提出的带 $p(u)=a^{2}u$ 的系统\eqref{align:ch05-relaxed-system}已广泛作为理解与阐释松弛特征的载体。我们将在后续章节考察该系统的理论，读者可在相应注释中找到相关文献。Baudin、Coquel 与 Tran [1] 提出上述松弛格式的一个变体，它与一维 Born--Infeld 系统有奇特关系；见 Serre [11]。我们也将在后续章节遇到松弛。

松弛与扩散之间的紧密关系还表现于带"摩擦"阻尼双曲系统解的长时间行为，特别在描述气体通过多孔介质等熵流动的简单系统中；见 Hsiao 与 Liu [1]、Tai-Ping Liu [25]、Serre 与 Xiao [1]、Hsiao 与 Luo [1]、Luo 与 Yang [1]、Nishihara 与 Yang [1]、Hsiao 与 Pan [2,3]、Hsiao、Li 与 Pan [1]、Hsiao 与 Li [1,2]、Nishihara、Wang 与 Yang [1,2,3]、Marcati 与 Mei [1]、He 与 Li [1]、Liu 与 Natalini [1]、Marcati 与 Pan [1]、Marcati 与 Nishihara [1]、Marcati、Mei 与 Rubino [1]、Pan [1,2]、Li 与 Saxton [1]、Huang 与 Pan [1,2]、Lattanzio 与 Rubino [1]、Huang、Marcati 与 Pan [1]、Di Francesco 与 Marcati [1]、Lan 与 Lin [1]、Dafermos 与 Pan [1]、Huang、Pan 与 Wang [1]。

在关于非双曲系统但表现出与带阻尼双曲系统相似行为的大量文献中，少量代表性样本如下。对 Euler--Poisson 系统，见 Poupaud、Rascle 与 Vila [1]、Dehua Wang [1,3]、Wang 与 Chen [1]、Guo [1]、Engelberg、Liu 与 Tadmor [1]、Li、Markowich 与 Mei [1]、Feldman、Ha 与 Slemrod [1]、Jang [1]、Chae 与 Tadmor [1] 与 Tadmor 与 Wei [1]。对半导体方程，见 Markowich、Ringhofer 与 Schmeiser 的专著 [1]，含完整参考文献列表；另见 Guo 与 Strauss [1]。对耦合 Euler 方程与计入辐射能通量的椭圆方程的辐射流体力学系统，见 Rohde 与 Yong [1]、Rohde、Wang 与 Xie [1] 与 Rohde 与 Xie [1,2]。Lions [2] 与 Feireisl [1] 的专著处理若干空间维数中可压缩粘弹性流体方程组，并提供详尽文献。当然，不可压缩情形（含经典 Navier--Stokes 方程）的文献浩瀚。粘性流体磁流体动力学在 Chen 与 Wang [4,5] 与 Dehua Wang [4] 中讨论。对辐射磁流体动力学方程，见 Rohde 与 Yong [2]。对一维热粘弹性系统，见 Dafermos 与 Hsiao [2] 与 Dafermos [12]。对一维热弹性方程，见 Slemrod [1]、Dafermos 与 Hsiao [3] 及 Racke 的综述文章 [1]。最后，对带衰减记忆诱导粘性的一维粘弹性方程，见 MacCamy [1]、Dafermos 与 Nohel [1]、Dafermos [15] 及 Renardy、Hrusa 与 Nohel 的专著 [1]。

初边值问题（含\cref{section:ch05-初边值问题}勾勒的材料细节）的详尽讨论见 Benzoni-Gavage 与 Serre [2]。
另见 Higdon 的综述文章 [1]。特别地，关于弹性动力学方程的\cref{theorem:ch05-clamped-elastic}取自 Dafermos 与 Hrusa [1]，而关于真空存在下 Euler 方程的\cref{theorem:ch05-vacuum}归功于 Coutand 与 Shkoller [2]。关于后一问题的相关结果，见 Coutand 与 Shkoller [1]、Makino [1]、Liu 与 Yang [1]、Liu、Xin 与 Yang [1]、Tong Yang [4]、Xu 与 Yang [1] 与 Jang 与 Masmoudi [1]。

关于稳定性问题的展望，见 Benzoni-Gavage、Rousset、Serre 与 Zumbrun [1]。另见 Benzoni-Gavage 与 Coulombel [1]。粘性消去法及所得边界层的性质与稳定性相关问题近年被积极研究；见 H.O. Kreiss [1]、Benabdallah 与 Serre [1]、Gisclon 与 Serre [1]、Gisclon [1]、Grenier 与 Gues [1]、Kreiss 与 Kreiss [1]、Xin [6]、Serre 与 Zumbrun [1]、Serre [14,17,24]、Joseph 与 LeFloch [1,2,3]、Rousset [1,2,3]、Métivier 与 Zumbrun [1,2] 及 Guès、Métivier、Williams 与 Zumbrun [5,6]。

\end{document}
